% Póster del TP final: overfitting por sesgo de selección (solo la estrategia long-short).
%
% Compilar con LuaLaTeX, desde esta carpeta:   latexmk -lualatex poster.tex
% Las figuras de figuras/ se generan, desde la carpeta del repo, con:   .venv/Scripts/python poster/figuras.py
% Todos los números salen de results/fa3e6998e744_full/ (los mismos de docs/resultados_LS.md).

\documentclass[20pt, a0paper, portrait, margin=0mm, innermargin=18mm,
               blockverticalspace=10mm, colspace=16mm, subcolspace=10mm]{tikzposter}

\usepackage{fontspec}
\usepackage{unicode-math}
\setmainfont{Lato}[Extension=.ttf, UprightFont=*-Regular, BoldFont=*-Bold,
                   ItalicFont=*-Italic, BoldItalicFont=*-BoldItalic]
\setsansfont{Lato}[Extension=.ttf, UprightFont=*-Regular, BoldFont=*-Bold,
                   ItalicFont=*-Italic, BoldItalicFont=*-BoldItalic]
\setmathfont{LeteSansMath.otf}  % matemática sin serifa, diseñada para acompañar a Lato
\usepackage[spanish, es-nolayout, es-noshorthands]{babel}
\usepackage{graphicx}
\usepackage{booktabs}
\usepackage{enumitem}
\usetikzlibrary{calc, patterns.meta}

% ---------------------------------------------------------------- Colores (los de las figuras)
\definecolor{azul}{HTML}{1C5CAB}        % hold-out
\definecolor{celeste}{HTML}{86B6EF}     % validación (CV)
\definecolor{azuloscuro}{HTML}{123F78}  % banda del título
\definecolor{tinta}{HTML}{1F1F1D}
\definecolor{gris}{HTML}{52514E}
\definecolor{grisclaro}{HTML}{C3C2B7}
\definecolor{fondo}{HTML}{EEF4FC}

\definecolorstyle{TP}{}{
    \colorlet{backgroundcolor}{white}
    \colorlet{framecolor}{white}
    \colorlet{titlefgcolor}{white}
    \colorlet{titlebgcolor}{azuloscuro}
    \colorlet{blocktitlebgcolor}{celeste}
    \colorlet{blocktitlefgcolor}{azul}
    \colorlet{blockbodybgcolor}{white}
    \colorlet{blockbodyfgcolor}{tinta}
    \colorlet{innerblocktitlebgcolor}{white}
    \colorlet{innerblocktitlefgcolor}{azul}
    \colorlet{innerblockbodybgcolor}{fondo}
    \colorlet{innerblockbodyfgcolor}{tinta}
    \colorlet{notefgcolor}{tinta}
    \colorlet{notebgcolor}{fondo}
    \colorlet{notefrcolor}{fondo}
}
\usecolorstyle{TP}
\usebackgroundstyle{Empty}

% ---------------------------------------------------------------- Título: banda de ancho completo
\definetitlestyle{Banda}{
    width=\paperwidth, roundedcorners=0, linewidth=0pt, innersep=18mm,
    titletotopverticalspace=0mm, titletoblockverticalspace=12mm, titlegraphictotitledistance=0pt
}{
    \fill[titlebgcolor] (\titleposleft,\titleposbottom) rectangle (\titleposright,\titlepostop);
    \fill[celeste] (\titleposleft,\titleposbottom) rectangle (\titleposright,\titleposbottom+5mm);
}
\usetitlestyle{Banda}

\makeatletter
\renewcommand\TP@maketitle{%
    \centering
    {\color{white}\bfseries\fontsize{80}{94}\selectfont \@title\par}
    \vspace{10mm}
    {\color{celeste}\fontsize{42}{52}\selectfont Identificación de la sobrestimación del Sharpe Ratio en estrategias
     long-short sin data leakage\par}
    \vspace{12mm}
    {\color{white}\bfseries\fontsize{38}{46}\selectfont \@author\par}
    \vspace{6mm}
    {\color{white!78!azuloscuro}\fontsize{30}{38}\selectfont \@institute\par}
    \vspace{4mm}
}
\makeatother

% ---------------------------------------------------------------- Bloques: título azul con una línea celeste debajo
\defineblockstyle{Limpio}{
    titlewidthscale=1, bodywidthscale=1, titleleft,
    titleoffsetx=0pt, titleoffsety=0pt, bodyoffsetx=0pt, bodyoffsety=0pt,
    bodyverticalshift=0pt, roundedcorners=0, linewidth=0pt,
    titleinnersep=8mm, bodyinnersep=8mm
}{
    \ifBlockHasTitle
        \fill[blocktitlebgcolor] ($(blocktitle.south west)+(\blocktitleinnersep,0)$)
            rectangle ($(blocktitle.south east)+(-\blocktitleinnersep,2.5mm)$);
    \fi
}
\useblockstyle{Limpio}

% ---------------------------------------------------------------- Texto
\setlength{\parindent}{0pt}
\setlength{\parskip}{0.45em}
% Texto alineado a la izquierda (en columnas angostas, el justificado abre demasiado los espacios). El cuerpo de los
% bloques va en un \parbox, así que se agrega \RaggedRight adentro de \block.
\usepackage{ragged2e}
\setlength{\RaggedRightParindent}{0pt}
% Espacios fijos alrededor de = y + en el texto (con \RaggedRight se estiraban), y sin cortes de línea después de ellos
\AtBeginDocument{\thickmuskip=5mu\relax \medmuskip=4mu\relax \relpenalty=10000 \binoppenalty=10000}
% Tamaño del texto de los bloques: tikzposter usa \large (24.9 pt); se agranda un poco para leerlo de más lejos
\newcommand{\cuerpo}{\fontsize{28.5}{34}\selectfont}
\makeatletter
\expandafter\patchcmd\csname\string\block\endcsname{\large\color{blockbodyfgcolor}}
    {\cuerpo\color{blockbodyfgcolor}\RaggedRight}{}{\PackageError{poster}{No se pudo parchear \string\block}{}}
\makeatother
\setlist[itemize]{leftmargin=1.1em, itemsep=0.35em, topsep=0.3em, parsep=0pt,
                  label={\color{azul}\textbullet}}

\newcommand{\SR}{\mathrm{SR}}
\newcommand{\SRhat}{\widehat{\mathrm{SR}}}
\newcommand{\destacado}[1]{%
    \par\vspace{0.3em}\noindent
    \colorbox{fondo}{\parbox{\dimexpr\linewidth-2\fboxsep\relax}{\RaggedRight#1}}\par\vspace{0.3em}}
\setlength{\fboxsep}{7mm}

% Figura de ancho completo de la columna y su explicación, un poco más chica que el cuerpo
\newcommand{\leyenda}{\fontsize{25}{30.5}\selectfont}
\newcommand{\figura}[2]{%
    \par\noindent\includegraphics[width=\linewidth]{#1}\par\vspace{3mm}
    {\leyenda\color{gris}#2\par}}

\tikzposterlatexaffectionproofoff  % sin el logo de tikzposter al pie

% ---------------------------------------------------------------- Diagramas
\tikzset{
    entrenamiento/.style={fill=celeste!30},
    test/.style={fill=celeste},
    purga/.style={pattern={Lines[angle=45, distance=2.4mm, line width=0.7mm]}, pattern color=gris!70},
    embargo/.style={fill=grisclaro!80},
    holdout/.style={fill=azul},
    muestra/.style={minimum width=9mm, minimum height=7mm, inner sep=0pt},
}

% Purged K-fold: una fila por split (6), con el fold de test, la purga, el embargo y el hold-out.
% Eje x en años desde enero de 2006; la purga y el borde se dibujan con h = 6.
\newcommand{\diagramacv}{%
\begin{tikzpicture}[x=11mm, y=1mm, font=\normalsize]
    \foreach \k in {0,...,5} {
        \pgfmathsetmacro{\y}{-11*\k}
        \pgfmathsetmacro{\ts}{2*\k}
        \pgfmathsetmacro{\te}{2*\k+2}
        \fill[entrenamiento] (0,\y) rectangle (12,\y-8);
        \ifnum\k>0 \fill[white] (\ts-0.5,\y) rectangle (\ts,\y-8); \fill[purga] (\ts-0.5,\y) rectangle (\ts,\y-8); \fi
        \ifnum\k<5
            \fill[embargo] (\te,\y) rectangle (\te+1,\y-8);
            \fill[white] (12-5/12,\y) rectangle (12,\y-8); \fill[purga] (12-5/12,\y) rectangle (12,\y-8);
        \fi
        \fill[test] (\ts,\y) rectangle (\te,\y-8);
    }
    \fill[holdout] (12.2,0) rectangle (20,-63);
    \node[white, align=center, font=\large] at (16.1,-31.5) {\bfseries Hold-out\\[1mm]2018--2025\\[3mm]se evalúa\\una sola vez};
    \foreach \a/\lab in {0/2006, 2/2008, 4/2010, 6/2012, 8/2014, 10/2016, 12.2/2018}
        \node[above, gris, inner sep=1.5mm] at (\a,0) {\lab};
    \draw[gris, line width=0.6mm] (0,-66) -- (0,-68.5) -- (12,-68.5) -- (12,-66);
    \node[below, gris, inner sep=2mm] at (6,-68.5) {Desarrollo 2006--2017: 144 meses, 6 splits};
    % Leyenda, en dos filas
    \foreach \x/\yy/\estilo/\texto in {0/-84/entrenamiento/Entrenamiento, 10/-84/test/Test (validación),
                                       0/-95/purga/Purga y borde ($h$ meses), 10/-95/embargo/Embargo (12 meses)} {
        \fill[white] (\x,\yy) rectangle (\x+0.8,\yy-7);
        \fill[\estilo] (\x,\yy) rectangle (\x+0.8,\yy-7);
        \node[right, gris] at (\x+0.9,\yy-3.5) {\texto};
    }
\end{tikzpicture}}

% CSCV: los 12 bloques anuales del desarrollo, 6 in-sample y 6 out-of-sample (una de las 924 particiones)
\newcommand{\diagramacscv}{%
\begin{tikzpicture}[x=1mm, y=1mm, font=\normalsize]
    \foreach \i/\is in {0/1, 1/0, 2/1, 3/1, 4/0, 5/0, 6/1, 7/0, 8/1, 9/0, 10/0, 11/1} {
        \pgfmathtruncatemacro{\anio}{6+\i}
        \ifnum\is=1
            \fill[celeste] (20*\i,0) rectangle (20*\i+17,-15);
        \else
            \draw[azul, line width=0.8mm, fill=white] (20*\i+0.4,-0.4) rectangle (20*\i+16.6,-14.6);
        \fi
        \node[text height=1.4ex, text depth=0pt] at (20*\i+8.5,-7.5) {\ifnum\anio<10 0\fi\anio};
    }
    \fill[celeste] (0,-22) rectangle (9,-29);
    \node[right, gris] at (10,-25.5) {IS: se elige la campeona};
    \draw[azul, line width=0.8mm, fill=white] (120.4,-22.4) rectangle (128.6,-28.6);
    \node[right, gris] at (130,-25.5) {OOS: se mide su Sharpe Ratio};
\end{tikzpicture}}

% ---------------------------------------------------------------- Contenido
\title{Sesgo de selección en estrategias de trading}
\author{Juan Cruz Mendoza, Juan Ignacio Catania y Camila Rodríguez}
\institute{Facultad de Ciencias Exactas y Naturales,
           Universidad de Buenos Aires \textperiodcentered{} 2026}

\begin{document}
\maketitle

\begin{columns}

% ================================================================ Columna 1
\column{0.333}

\block{1. Introducción}{
Al desarrollar una estrategia en el mundo de las finanzas cuantitativas, es muy habitual probar muchas
variantes de la misma idea y quedarse con la de mejor backtest. Aunque cada configuración se valide sin leakage, el
máximo de muchas estimaciones ruidosas tiene un \textbf{sesgo de selección}, y la elegida suele
rendir mucho peor out-of-sample. 

\textbf{Objetivo:} mostrar lo fácil que es sobreajustar por este sesgo y evaluar si
el DSR, la CSCV y un hold-out permiten detectarlo y medirlo.
}

\block{2.1 \,Datos y estrategia}{
\begin{itemize}
    \item \textbf{Universo:} las 40 empresas más grandes del S\&P~500 en 2005
          que siguieron cotizando hasta 2025, con precios diarios.
    \item \textbf{Features}: momentum 12-1, momentum 3 meses, reversión 1 mes, volatilidad
          6 meses y tamaño (volumen en dólares de 12 meses).
    \item \textbf{Target:} 1 si el retorno de la acción en los próximos $h$ meses está en el cuantil superior del
          mes (tercil o quintil).
    \item \textbf{Modelo base:} Random Forest de 200 árboles.
    \item \textbf{Estrategia long-short:} cada $h$ meses se compran las acciones con mayor probabilidad predicha y
          se hace un short con las de menor probabilidad, con pesos iguales.
    \item \textbf{Costos:} 10bps por lado sobre el turnover de cada rebalanceo.
    \item \textbf{Métrica:} Sharpe Ratio neto de costos, anualizado.
\end{itemize}

\vspace{0.4em}
% Sin \extracolsep{\fill}: dentro de tikzposter, \fill es el comando de TikZ
{\fontsize{22}{26.5}\selectfont
\begin{tabular*}{\linewidth}{@{\extracolsep{0pt plus 1fil}}llr@{}}
    \toprule
    \textbf{Dimensión} & \textbf{Valores} & \textbf{Opciones} \\
    \midrule
    Subconjunto de features & todas las combinaciones & 31 \\
    Horizonte $h$ & 1, 3 y 6 meses & 3 \\
    Cuantil del target & tercil, quintil & 2 \\
    Profundidad máxima & 2, 4 y 8 & 3 \\
    Mínimo de muestras por hoja & 1 y 50 & 2 \\
    \midrule
    \textbf{Configuraciones} & $31 \times 3 \times 2 \times 3 \times 2$ & \textbf{1116} \\
    \bottomrule
\end{tabular*}\par}

\vspace{0.3em}
Se elige la configuración de mayor Sharpe Ratio en la validación cruzada (CV).
}

\block{2.2 \,Validación sin leakage}{
\diagramacv

\vspace{0.2em}
\begin{itemize}
    \item \textbf{Purged K-fold} ($K = 6$) en el desarrollo: una serie out-of-sample de 144 meses por configuración.
    \item \textbf{Purga:} sin muestras cuya etiqueta se solape con el test o con el hold-out.
    \item \textbf{Embargo:} 12 meses después del test, porque las features miran 12 meses hacia atrás.
    \item \textbf{Hold-out:} la elegida se evalúa una sola vez, al final (las 1116 también, solo para medir).
\end{itemize}
}

% ================================================================ Columna 2
\column{0.333}

\block{2.3 \,Deflated Sharpe Ratio (DSR)}{
Si las $N$ configuraciones fueran independientes y ninguna tuviera señal, sus Sharpe Ratio serían
aproximadamente normales, $\SRhat_n \sim \mathcal{N}(0, \sigma^2)$, y su máximo seguiría 
aproximadamente una distribución de Gumbel. Su esperanza es el máximo esperado por azar:
\[
    \SR_0 = \sigma \left[ (1-\gamma)\, \Phi^{-1}\!\left(1-\frac{1}{N}\right)
            + \gamma\, \Phi^{-1}\!\left(1-\frac{1}{N e}\right) \right],
\]
El DSR es la probabilidad de que el Sharpe Ratio
verdadero de la elegida supere a $\SR_0$, corrigiendo por la asimetría, la curtosis y el largo de la serie. Como las
configuraciones están correlacionadas ($\bar\rho = 0.58$), se usa el $N$ efectivo
$N_{\mathrm{ef}} = \bar\rho + (1-\bar\rho)\, N$, que da $\SR_0 = 0.70$. La elegida es
significativa si $\mathrm{DSR} > 0.95$.
}

\block{2.4 \,CSCV}{
La Combinatorially Symmetric Cross-Validation parte los 144 meses de retornos out-of-sample de las 1116
configuraciones en 12 bloques anuales:

\vspace{0.2em}
\diagramacscv

En cada una de las formas de elegir 6 bloques in-sample (IS), se toma la configuración de
mayor Sharpe Ratio IS (la campeona) y se mide su Sharpe Ratio en los otros 6 bloques (OOS). El promedio sobre las
particiones es la \textbf{estimación CSCV}: estima lo que rinde out-of-sample el procedimiento de elegir la
mejor.
}

\block{3.1 \,La ventaja de la elegida no supera al máximo esperado por azar}{
\figura{figuras/fig1_max_sharpe.pdf}{Sharpe Ratio CV máximo entre $n$ configuraciones sorteadas (promedio de 1000
sorteos). Se mide si la elegida supera en SR al resto.}
}

% ================================================================ Columna 3
\column{0.333}

\block{3.2 \,En el hold-out, la elegida cae y el promedio no}{
\figura{figuras/fig2_elegida_vs_configuraciones.pdf}{La elegida cae de 0.96 a 0.53 y el promedio no se mueve: la
caída se debe a la selección. La estimación CSCV (0.40) anticipó la caída.}
}

\block{3.3 \,La CV ordena, pero sobrestima a las mejores}{
\figura{figuras/fig3_deciles.pdf}{Decil 1 es el 10\,\% más bajo, decil 10 es el
10\,\% más alto. En las puntas, los valores vuelven hacia el promedio: es regresión a la media, porque en los
extremos pesa más el ruido.}
}

\block{3.4 \,Elegir por AUC no evita el problema}{
\figura{figuras/fig4_sharpe_vs_auc.pdf}{En cada grupo $(h, q)$ se elige la mejor por Sharpe Ratio y la mejor por AUC
de la CV: el gráfico muestra el promedio de los 6 grupos. Elegir por AUC evita el sesgo optimista, pero rinde menos.}
}

\end{columns}

\begin{columns}
\column{0.5}
\block{3.5 \,Resumen}{
\leyenda  % la tabla, al tamaño de las leyendas: con el del cuerpo, los encabezados quedan pegados
\begin{tabular*}{\linewidth}{@{\extracolsep{0pt plus 1fil}}lccc@{}}
    \toprule
    & \textbf{Sharpe Ratio CV} & \textbf{Sharpe Ratio hold-out} & \textbf{DSR} \\
    & (2006--2017) & (2018--2025) & (CV) \\
    \midrule
    RF elegida por Sharpe Ratio & 0.96 & 0.53 & 0.80 \\
    RF elegida por AUC & 0.56 & 0.35 & -- \\
    Promedio de las 1116 configuraciones & 0.17 & 0.17 & -- \\
    \bottomrule
\end{tabular*}

\vspace{0.6em}
{\leyenda\color{gris}La elegida por AUC es la de mayor AUC de la CV (grupo $h = 6$, $q = 5$, el mismo de la elegida
por Sharpe Ratio).\par}

% Referencias debajo de la tabla, en el espacio que deja la columna de las conclusiones (abajo no entran)
\vspace{1.2em}
{\small\color{gris}
\textbf{Referencias.}
Bailey, D. H. y López de Prado, M. (2014). The Deflated Sharpe Ratio: correcting for selection bias, backtest
overfitting and non-normality. \textit{The Journal of Portfolio Management}, 40(5), 94--107. \quad
Bailey, D. H., Borwein, J., López de Prado, M. y Zhu, Q. J. (2017). The probability of backtest overfitting.
\textit{Journal of Computational Finance}, 20(4), 39--69. \quad
López de Prado, M. (2018). \textit{Advances in Financial Machine Learning}. Wiley.\par}

% El link al código, en su propia línea y más grande que las referencias
\vspace{0.8em}
{\leyenda\textbf{Código:} {\color{azul}github.com/JuanCruzMendoza/quant-selection-bias}\par}
}

\column{0.5}
\block{4. Conclusiones}{
\begin{itemize}
    \item \textbf{Sobreajustar por selección es fácil.} Con purged K-fold, embargo y un hold-out intacto, el
          Sharpe Ratio de validación de la elegida (0.96) casi duplicó al del hold-out (0.53). Evitar el leakage no
          alcanza: el sesgo crece con la cantidad de intentos.
    \item \textbf{Las herramientas lo detectan, si se registra todo.} El DSR y la estimación CSCV anticiparon la
          caída usando solo la data de desarrollo, pero necesitan conocer todos los intentos y sus retornos. Con una
          metodología rigurosa que guarda todas las corridas, podemos saber cuánto se “quemaron” los datos.
    \item \textbf{Limitación:} survivorship bias. Las 40 acciones cotizaron durante todo 2005--2025,
          así que quedan afuera las empresas que quebraron o fueron adquiridas.
\end{itemize}
}
\end{columns}

\end{document}
